\hypertarget{the-capability-mesh}{%
\chapter{11. The capability mesh}\label{the-capability-mesh}}

By the time a kennel is running, the design has spent four chapters
making sure it cannot see any other. Its filesystem holds no other
kennel's files; its network reaches no other kennel's services; its
local surface names no other kennel's sockets. Two kennels on the same
machine are, by default, as invisible to each other as two machines
(\texttt{construction-by-absence}). That is usually right, but not
always. Real work sometimes needs one confined workload to use a
capability another confined workload provides: an application that must
render to a display server it does not itself run, reach a session bus
through a broker it does not host, read from a cache that lives in a
kennel of its own. The mesh is the deliberate exception that allows
this: without either kennel gaining the other's authority, and without
disturbing the rule that kennels have no hold on one another.

The exception is kept as narrow as the rule. A provider and a consumer
remain siblings under the framework, neither nested in the other; the
only thing that passes between them is the single connector the operator
declared and the framework brokered. All this adds is how that one seam
is opened safely between two confinements that still cannot, in any
other respect, see each other.

\hypertarget{name-not-peer}{%
\section{11.1 Name, not peer}\label{name-not-peer}}

A capability enters the mesh by being offered and reached, and the two
declarations are deliberately blind to each other. A provider offers a
capability under a public name (the way a service announces an
interface, not the way a process announces itself) and lists nothing
about who may consume it. A consumer reaches for a capability by that
same name, and names no provider. Nothing in either declaration points
at another kennel; there is no peer reference anywhere on the policy
surface. A consumer asks for the display, or the cache, not for a
kennel, and the framework resolves that name, at runtime, to whoever
offers it.

That blindness is what lets the policies stay independent. Because no
declaration names another kennel, each policy is complete on its own (it
can be written, compiled, and signed with no other kennel in view), and
the cross-kennel reach it asks for is resolved later, against whatever
providers the operator has installed. The reach is real, but it is never
authored as a link; it is two independent declarations that the
framework, holding both, recognises as a matched pair.

And the matching is deny-by-default, the discipline carried straight
from delegation. A workload cannot author a cross-kennel reach any more
than it could author a policy (\texttt{request-dont-author}); both the
offer and the use are signed declarations, and the framework brokers
only the pairs where both sides declared and the operator signed both.
What a kennel can reach is a projection of those declarations, computed
fresh, never a standing list of grants left lying in the daemon. Where a
public name is not enough (where more than one kennel might offer the
same name and a consumer must be sure of the one it meant), provider and
consumer can pair on a second key, a private one this time. It is not a
secret in the sense of a credential whose strength is its concealment;
it is the other half of the lookup. The public name takes a consumer to
everything offered under it, the way a title takes you to a shelf; the
private key picks the one intended provider from among them, the way a
catalogue number picks the one volume. Both are lookup keys (one
advertised, one not), and the framework requires both to match before it
brokers the pair. The private key names no kennel either: it is a value
the intended provider and consumer happen to share so the match lands on
the right one, not a reference either holds to the other.

\hypertarget{resolution-is-runtime}{%
\section{11.2 Resolution is runtime}\label{resolution-is-runtime}}

A cross-kennel match cannot be settled when a policy is signed, and the
design is careful not to pretend it can. Each policy is compiled and
signed alone (the compiler only ever holds the one in front of it), and
the providers that can answer a consumer are whatever the operator has
installed by the time the consumer runs, not whatever happened to exist
when the consumer was signed. To promise at signing time that a given
reach will resolve would be to make a signed, immutable artefact depend
on a separate, mutable, independently-authored one: a promise the system
has no way to keep. So matching is a runtime property, and the design
puts it there rather than faking it earlier.

What can be checked early is checked early, and only that: holding one
policy, the compiler confirms each offer and each use is well-formed and
that a reserved name is one the policy is entitled to claim: everything
decidable from the single artefact, nothing that consults another
kennel. The rest waits for construction, when the consumer's name is
resolved against the catalogue of what is installed. A consumer says, of
each capability it reaches, whether it can live without it. A required
capability that resolves to nothing (no installed provider offers the
name) fails the kennel's construction loudly, rather than starting a
workload missing something it depends on; an optional one that resolves
to nothing is simply skipped, and the workload bears the absence. It is
the hard-versus-soft dependency distinction systemd draws between what a
unit requires and what it merely wants, evaluated here against what the
deployment has rather than what a policy hoped for.

\hypertarget{typed-shapes}{%
\section{11.3 Typed shapes}\label{typed-shapes}}

A capability also carries its shape (the kind of transport by which it
is reached), and the consumer must state the shape it expects, because
the ways of reaching a capability are not interchangeable. A connected
socket, a name made reachable on the bus, and a brokered channel are
acquired through entirely different means; a consumer that expects one
and is handed another has been misdelivered, not served. So the broker
matches shape as well as name, and refuses a mismatch at the moment of
connection rather than letting a consumer find it as a failure later.

The set of shapes is small and closed on purpose, and the reason is the
discipline that has governed every facade so far. Each shape is a
transport the framework already mediates: the local socket of the
services chapter, the bus reached through that same chapter's broker,
the channel that carries a delegated sibling. The mesh adds no new way
to move bytes between kennels; it adds a way to declare and broker the
ways that already exist. A single catalogue and a single broker describe
every cross-kennel reach uniformly, and widening the mesh means teaching
it another transport the framework already brokers, never a new wire
protocol in the trusted path (\texttt{do-less}). The broker itself stays
where every mediator in this design stays: it matches, it authorises, it
mints the connector, and it steps out of the bytes that then flow
(\texttt{control-not-data-plane}).

One thing the broker never does is hand a consumer a connection it has
not yet asked for. A capability is reached on demand, when the workload
acts, not pressed into the kennel at construction, which is also what
lets a provider stay down until it is first reached and come up only
then, so a standing service that no running kennel is currently
consuming costs nothing to have offered.

\hypertarget{the-standing-provider}{%
\section{11.4 The standing provider}\label{the-standing-provider}}

What the mesh carries that delegation did not is permanence. A delegated
sibling was ephemeral: spun up for a task, reaped when the requester was
done. A mesh provider is usually a standing service: a confined kennel
the operator enables once and leaves up, serving many consumers across
its life: the confined display service the earlier chapters pointed
toward, a session-bus broker, a shared cache. Permanence changes the
trust the provider carries and the residual it leaves (\texttt{T3.10}).

A provider is, almost by definition, code that needs a foothold in the
host: the display service must reach the real display, the bus broker
the real bus, the cache the real storage. The ordinary way to give code
a foothold in the host is to run it in the host: unconfined, at the
user's full authority, trusted because there seemed nowhere else to put
it. Kennel puts it in a kennel instead. A provider is itself a confined
workload, granted exactly the one leg into the host its job requires and
nothing more: the same confinement-and-enablement structure every other
kennel runs under, turned on the very services that would otherwise be
the unconfined host-side code everything else depends on. The trust it
carries is layered accordingly, by the mechanism delegation already
established. A reserved name cannot be authored by a policy that simply
fancies it: a reserved namespace belongs to a tier, and only that tier's
key may sign a template claiming a name within it. The
\texttt{org.projectkennel.*} namespace is the project maintainers',
always: a name under it is honoured only on a maintainer-signed
template, and no host or operator signature can claim one. A host may
reserve a namespace of its own in turn (\texttt{com.acme.*}, say) for
the standing services it runs on its own machines; a name under that
prefix is honoured only on a template signed by that host's key. The
prefix is the claim of authority and the signature is its proof, and the
two must agree or the name is refused. What the operator does is take
such a signed template (the maintainers' or the host's) and adapt their
own version of it, filling the fields it leaves mutable, inheriting the
reserved name and the rest exactly as signed, and enable it, the
enablement being the grant of the host foothold. The vouching is layered
to match: the tier that owns the namespace signs the reserved shape, and
the operator stands up this instance of it. The shared thing that many
confinements lean on is held to a higher bar than the private thing one
workload holds: signed, at its origin, by the tier that owns the name,
and bounded, where the operator varies it, to the blanks its author left
open.

The residual is the other face of the same permanence. A standing
provider that is compromised is compromised for every consumer that
reaches it, for as long as it stays up: a wider blast than an ephemeral
sibling's, whose damage ends with its task. The mesh does not erase that
exposure; it replaces a worse one. Without it, the display, the bus, the
cache are unconfined host code, and an application reaches them at the
user's full authority: a compromise anywhere in that chain is host-wide.
With it, those services are themselves confined to the one host leg each
was granted, and the application reaches them through a declared,
brokered, deny-by-default seam; a compromise is bounded to a provider
that holds little, and to the consumers that declared they would trust
it. The exposure is named, narrowed, and in view: far smaller than the
ambient reach it replaces, and never pretended away.

With the mesh, the resource classes are complete. Every kind of thing a
workload touches (its files, what it may run, the network, the local
services of its session, the images it runs inside, the siblings it
spawns, and now the capabilities it shares with other confinements) has
been answered by the same two limbs and the same handful of principles.
The mesh comes last because it is the one exception to the isolation all
the others build, and it earns the exception by holding to the rules
they established: declared, not authored; brokered, not ambient;
deny-by-default; and narrowed to the single seam the operator opened.
What remains, in the chapters that follow, is not another resource but
the grant itself: how a policy is written, composed, signed, and made to
fail safe, the declared authority every one of these chapters has taken
as given.
